\section{Proofs}\label{app:proofs}

Throughout, $\|\cdot\|_1$ is the $\ell_1$ norm and constants $c,C$ may change from line to line.

\paragraph{Covering the simplex.}
For $\epsilon\in(0,1]$ there is a set $\mathcal N_\epsilon\subset\Delta_M$ with $|\mathcal N_\epsilon|\le(3/\epsilon)^M$ such that every $\bm w\in\Delta_M$ has a point $\bm w'\in\mathcal N_\epsilon$ with $\|\bm w-\bm w'\|_1\le\epsilon$ (volumetric argument for the $\ell_1$ ball, e.g.\ \citealp[Ch.~4]{Vershynin2018}).
Hence $\mathcal N_\epsilon^G$ is an $\epsilon$-net of $\Delta_M^G$ in the norm $\max_g\|\bm w_g-\bm w'_g\|_1$, of cardinality at most $(3/\epsilon)^{GM}$.

\paragraph{Lipschitz property.}
For all $(\bm\gamma,\bm W)$ and $(\bm\gamma,\bm W')$ with $\max_g\|\bm w_g-\bm w_g'\|_1\le\epsilon$, Assumption~\ref{as:bounded} gives
$|\widehat R(\bm\gamma,\bm W)-\widehat R(\bm\gamma,\bm W')|\le\bar aB\epsilon$ and $|R(\bm\gamma,\bm W)-R(\bm\gamma,\bm W')|\le\bar aB\epsilon$.

\subsection*{Proof of Theorem~\ref{thm:oracle}}

Fix $\epsilon=1/(NT_e)$.
The set $\mathcal F=\{1,\dots,G\}^N\times\mathcal N_\epsilon^G$ has cardinality at most $G^N(3NT_e)^{GM}$.
By Assumption~\ref{as:concentration} and the union bound,
\[
\Pr\Big(\max_{(\bm\gamma,\bm W)\in\mathcal F}|\widehat R(\bm\gamma,\bm W)-R(\bm\gamma,\bm W)|>x\Big)\le 2G^N(3NT_e)^{GM}\exp(-cNT_ex^2/L^2).
\]
Setting the right-hand side equal to $\delta$ gives, with probability at least $1-\delta$,
$\max_{\mathcal F}|\widehat R-R|\le x_\delta:=L\{(GM\log(3NT_e)+N\log G+\log(2/\delta))/(cNT_e)\}^{1/2}$.
For an arbitrary $(\bm\gamma,\bm W)\in\mathcal P_G$ let $\bm W'$ be its nearest net point; by the Lipschitz property,
\[
\sup_{\mathcal P_G}|\widehat R-R|\le x_\delta+2\bar aB\epsilon .
\]
Let $(\bm\gamma^*,\bm W^*)$ satisfy $R(\bm\gamma^*,\bm W^*)\le R^*_G+\eta$ for arbitrary $\eta>0$.
Since $(\widehat{\bm\gamma},\widehat{\bm W})$ minimises $\widehat R$,
\[
R(\widehat{\bm\gamma},\widehat{\bm W})\le\widehat R(\widehat{\bm\gamma},\widehat{\bm W})+x_\delta+2\bar aB\epsilon
\le\widehat R(\bm\gamma^*,\bm W^*)+x_\delta+2\bar aB\epsilon\le R^*_G+\eta+2x_\delta+4\bar aB\epsilon .
\]
Letting $\eta\downarrow0$ and substituting $\epsilon=1/(NT_e)$ gives~\eqref{eq:oracle}. \qed

\subsection*{Proof of Corollary~\ref{cor:ao}}

Theorem~\ref{thm:oracle} with $\delta=\delta_{N,T}\to0$ slowly (for instance $\log(1/\delta)=\log(NT_e)$) shows $R(\widehat{\bm\gamma},\widehat{\bm W})-R^*_G=O_p(r_{N,T})$ with
$r_{N,T}=\{GM\log(NT_e)/(NT_e)+\log G/T_e+\log(NT_e)/(NT_e)\}^{1/2}\to0$.
Dividing by $R^*_G\ge\kappa$ yields $R(\widehat{\bm\gamma},\widehat{\bm W})/R^*_G-1=O_p(r_{N,T}/\kappa)\to_p0$.
Because $R(\widehat{\bm\gamma},\widehat{\bm W})\ge R^*_N$ always and $R^*_G\ge R^*_N$, the second claim follows once we show $R^*_G=R^*_N$ under the stated structure.
Let $\bm w^0_g$ be a common minimiser of $R_i$ for all $i$ with $\gamma^0_i=g$ (padding with arbitrary vectors if $G_0<G$); then
$R(\bm\gamma^0,\bm W^0)=N^{-1}\sum_i\min_{\bm w}R_i(\bm w)=R^*_N$, so $R^*_G\le R^*_N$, and the reverse inequality holds because $\mathcal P_G$ is a subset of the unit-specific class. \qed

\subsection*{Proof of Theorem~\ref{thm:recovery}}

\emph{Step 1 (average consistency).}
Under Assumption~\ref{as:groups}(a), $R^*_G=R(\bm\gamma^0,\bm W^0)=R^*_N$ and
\[
\frac\lambda N\sum_{i=1}^N\big\|\widehat{\bm w}_{\widehat\gamma_i}-\bm w^0_{\gamma^0_i}\big\|_2^2\le R(\widehat{\bm\gamma},\widehat{\bm W})-R^*_G=o_p(1)
\]
by Theorem~\ref{thm:oracle}, since $\log N/T_e\to0$ implies $\log G/T_e\to0$ and $GM\log(NT_e)/(NT_e)\to0$ for fixed $G$ and $M$.

\emph{Step 2 (consistency of the segment weights).}
Fix $\eta\in(0,d/4)$.
By Step~1 and Markov's inequality, the fraction of units with $\|\widehat{\bm w}_{\widehat\gamma_i}-\bm w^0_{\gamma^0_i}\|_2>\eta$ is $o_p(1)$.
Since every true segment $g$ contains at least $\pi N$ units, with probability approaching one each $g$ contains a unit $i$ with $\|\widehat{\bm w}_{\widehat\gamma_i}-\bm w^0_g\|_2\le\eta$; call $h(g)=\widehat\gamma_i$.
If $h(g)=h(g')$ for $g\ne g'$, the triangle inequality gives $\|\bm w^0_g-\bm w^0_{g'}\|_2\le2\eta<d$, contradicting Assumption~\ref{as:groups}(b).
Thus $h$ is injective, hence a bijection of $\{1,\dots,G\}$, and after relabelling $\max_g\|\widehat{\bm w}_g-\bm w^0_g\|_2\le\eta$ with probability approaching one.
Since $\eta$ is arbitrary, (i) follows.

\emph{Step 3 (memberships).}
On the event $E_\eta=\{\max_g\|\widehat{\bm w}_g-\bm w^0_g\|_2\le\eta\}$, unit $i$ with $\gamma^0_i=g$ can be assigned to $h\ne g$ only if $\widehat R_i(\widehat{\bm w}_h)\le\widehat R_i(\widehat{\bm w}_g)$; every global minimiser assigns each unit to a segment minimising its own empirical cost, since otherwise reassigning the unit would lower~\eqref{eq:gdma}.
By Assumption~\ref{as:groups} and the Lipschitz property (with $\|\cdot\|_1\le\sqrt M\|\cdot\|_2$),
\[
R_i(\widehat{\bm w}_h)-R_i(\widehat{\bm w}_g)\ge\lambda(d-\eta)^2-\bar aB\sqrt M\eta\ge\lambda d^2/4=:\Delta
\]
for $\eta$ small enough.
Hence misclassification of unit $i$ on $E_\eta$ requires $\sup_{\bm w\in\Delta_M}|\widehat R_i(\bm w)-R_i(\bm w)|\ge\Delta/2$.
Covering $\Delta_M$ by $\mathcal N_\epsilon$ with $\epsilon=\Delta/(8\bar aB)$ and applying the unit-level part of Assumption~\ref{as:concentration},
\[
\Pr\Big(\sup_{\bm w}|\widehat R_i(\bm w)-R_i(\bm w)|\ge\Delta/2\Big)\le2\Big(\frac{24\bar aB}{\Delta}\Big)^M\exp\!\Big(-\frac{cT_e\Delta^2}{16L^2}\Big)=:C_M e^{-c'T_e}.
\]
Because the supremum is over all $\bm w$, the bound does not require independence between $\widehat{\bm W}$ and unit $i$'s data.
A union bound over units gives $\Pr(\widehat{\bm\gamma}\ne\bm\gamma^0)\le\Pr(E_\eta^c)+NC_Me^{-c'T_e}$, which tends to zero because $\log N/T_e\to0$. This is (ii).

\emph{Step 4 (oracle equivalence).}
On the event $\{\widehat{\bm\gamma}=\bm\gamma^0\}$, $\widehat{\bm W}$ minimises $\widehat R(\bm\gamma^0,\cdot)$, so $\widehat{\bm W}=\widetilde{\bm W}$ (if the minimiser is not unique, the two estimators select from the same set).
By (ii) this event has probability approaching one. \qed

\section{Out-of-sample risk}\label{app:oos}

The theory concerns the population risk $R$ evaluated at the estimated policy.
In a rolling application the policy estimated on periods $1,\dots,T$ is applied in periods $T+1,\dots,T+k$.
If $\{\bm z_{it}\}$ is $\beta$-mixing with coefficients $\beta(\cdot)$, a coupling argument \citep{Yu1994} shows that the expected cost of the estimated policy in period $T+s$ differs from $\mathbb E\,R(\widehat{\bm\gamma},\widehat{\bm W})$ by at most $L\beta(s)$.
The bounds of Section~\ref{sec:theory} therefore also bound the out-of-sample cost up to a term that vanishes with the gap between estimation and application.

\section{Balancing authorities}\label{app:data}

\begin{table}[h]
\centering
\caption{The 46 balancing authorities, ordered by mean hourly demand over 2019--2026. Flex.: interchange flexibility in 2019 (range between the 5th and 95th percentiles of hourly total interchange, divided by mean demand). $\tau_i$: calibrated critical ratio.}\label{tab:bas}
\scriptsize
\begin{tabular}{@{}llrrr@{}}
\toprule
Code & Balancing authority & Mean MW & Flex. & $\tau_i$ \\
\midrule
PJM & PJM Interconnection, LLC & 92,507 & 0.06 & 0.947 \\
MISO & Midcontinent Independent System Operator, Inc. & 74,189 & 0.07 & 0.943 \\
ERCO & Electric Reliability Council of Texas, Inc. & 49,750 & 0.03 & 0.950 \\
SWPP & Southwest Power Pool & 32,259 & 0.08 & 0.940 \\
SOCO & Southern Company Services, Inc. - Trans & 27,075 & 0.12 & 0.927 \\
CISO & California Independent System Operator & 25,319 & 0.29 & 0.900 \\
TVA & Tennessee Valley Authority & 18,637 & 0.09 & 0.937 \\
NYIS & New York Independent System Operator & 17,306 & 0.18 & 0.910 \\
FPL & Florida Power \& Light Co. & 15,873 & 0.10 & 0.933 \\
ISNE & ISO New England & 13,228 & 0.16 & 0.920 \\
DUK & Duke Energy Carolinas & 12,212 & 0.23 & 0.907 \\
CPLE & Duke Energy Progress East & 7,045 & 0.37 & 0.880 \\
BPAT & Bonneville Power Administration & 6,631 & 0.84 & 0.827 \\
PACE & PacifiCorp East & 5,803 & 0.42 & 0.877 \\
PSCO & Public Service Company of Colorado & 5,312 & 0.11 & 0.930 \\
NEVP & Nevada Power Company & 4,523 & 0.33 & 0.883 \\
LGEE & Louisville Gas and Electric Company and Kentucky Utilities Company & 4,166 & 0.17 & 0.913 \\
AZPS & Arizona Public Service Company & 3,893 & 0.67 & 0.843 \\
SC & South Carolina Public Service Authority & 2,901 & 0.32 & 0.890 \\
LDWP & Los Angeles Department of Water and Power & 2,867 & 0.30 & 0.897 \\
SCEG & Dominion Energy South Carolina, Inc. & 2,814 & 0.15 & 0.923 \\
AECI & Associated Electric Cooperative, Inc. & 2,658 & 0.52 & 0.867 \\
TEC & Tampa Electric Company & 2,530 & 0.26 & 0.903 \\
PGE & Portland General Electric Company & 2,523 & 0.53 & 0.863 \\
PACW & PacifiCorp West & 2,392 & 0.55 & 0.860 \\
FMPP & Florida Municipal Power Pool & 2,135 & 0.17 & 0.917 \\
IPCO & Idaho Power Company & 2,097 & 0.79 & 0.833 \\
BANC & Balancing Authority of Northern California & 1,993 & 0.59 & 0.853 \\
PNM & Public Service Company of New Mexico & 1,645 & 0.64 & 0.847 \\
TEPC & Tucson Electric Power & 1,591 & 0.71 & 0.837 \\
JEA & JEA & 1,532 & 0.32 & 0.893 \\
AVA & Avista Corporation & 1,460 & 0.70 & 0.840 \\
NWMT & NorthWestern Corporation & 1,358 & 1.07 & 0.810 \\
SCL & Seattle City Light & 1,076 & 0.92 & 0.817 \\
WALC & Western Area Power Administration - Desert Southwest Region & 1,045 & 0.95 & 0.813 \\
EPE & El Paso Electric Company & 1,029 & 0.46 & 0.873 \\
GCPD & Public Utility District No. 2 of Grant County, Washington & 696 & 1.58 & 0.807 \\
CPLW & Duke Energy Progress West & 564 & 0.61 & 0.850 \\
TPWR & City of Tacoma, Department of Public Utilities, Light Division & 540 & 0.56 & 0.857 \\
IID & Imperial Irrigation District & 441 & 0.90 & 0.820 \\
TIDC & Turlock Irrigation District & 328 & 0.81 & 0.830 \\
TAL & City of Tallahassee & 322 & 0.32 & 0.887 \\
DOPD & PUD No. 1 of Douglas County & 245 & 2.21 & 0.803 \\
GVL & Gainesville Regional Utilities & 240 & 0.50 & 0.870 \\
CHPD & Public Utility District No. 1 of Chelan County & 218 & 4.71 & 0.800 \\
HST & City of Homestead & 71 & 0.89 & 0.823 \\
\bottomrule
\end{tabular}
\end{table}

\section{Additional results for the global panel}\label{app:global}

\begin{table}[h]
\centering
\caption{Global panel with the main candidate set (eleven candidates including the weather models): system cost relative to the reference forecaster for estimation windows of 14 and 28 days (world).}\label{tab:globalwindows}
\small
\begin{tabular}{@{}lcc@{}}
\toprule
Method & $W=14$ & $W=28$ \\
\midrule
Reference forecaster$^a$ & 1.000 & 1.000 \\
Equal weights & 1.217 & 1.217 \\
Best single (per area) & 0.943 & 0.923 \\
Pooled DF weights & 0.906 & 0.902 \\
Per-area DF weights & 0.918 & 0.895 \\
Shrinkage to pooled & 0.905 & 0.888 \\
Two-step $k$-means & 0.899 & 0.884 \\
Quantile-regression comb. & 1.174 & 1.016 \\
Forecast-then-commit, per area & 1.030 & 0.988 \\
Forecast-then-commit, grouped & 0.990 & 0.966 \\
Neural gate (deep MoE) & -- & 0.956 \\
GDMA & 0.900 & 0.884 \\
\textbf{Soft GDMA} & 0.902 & 0.884 \\
Fair soft GDMA & 0.897 & 0.881 \\
\bottomrule
\end{tabular}
\end{table}

\begin{table}[h]
\centering
\caption{Global panel with the main candidate set and with a common critical ratio $\tau=0.9$ for every grid area ($W=28$): system cost relative to the reference forecaster. The neural gate was not run for this specification.}\label{tab:globaltau}
\small
\begin{tabular}{@{}lcccccc@{}}
\toprule
Method & World & North America & Europe & Oceania & Asia & Africa \\
\midrule
Reference forecaster$^a$ & 1.000 & 1.000 & 1.000 & 1.000 & 1.000 & 1.000 \\
Equal weights & 1.193 & 1.401 & 1.059 & 0.996 & 1.035 & 0.976 \\
Best single (per area) & 0.917 & 0.854 & 0.942 & 1.008 & 1.023 & 0.915 \\
Pooled DF weights & 0.899 & 0.863 & 0.914 & 0.983 & \textbf{0.935} & 0.926 \\
Per-area DF weights & 0.891 & 0.832 & 0.912 & 0.989 & 0.990 & 0.932 \\
Shrinkage to pooled & 0.884 & 0.830 & 0.905 & 0.984 & 0.971 & 0.925 \\
Two-step $k$-means & 0.879 & 0.833 & 0.896 & 0.975 & 0.946 & 0.895 \\
Quantile-regression comb. & 0.987 & 0.945 & 0.987 & 0.983 & 1.179 & 1.064 \\
Forecast-then-commit, per area & 0.958 & 0.934 & 0.932 & 0.989 & 1.184 & 0.946 \\
Forecast-then-commit, grouped & 0.938 & 0.932 & 0.898 & 0.981 & 1.145 & 0.926 \\
GDMA & 0.878 & 0.827 & 0.899 & \textbf{0.973} & 0.954 & \textbf{0.891} \\
\textbf{Soft GDMA} & 0.880 & 0.824 & 0.901 & 0.976 & 0.972 & 0.908 \\
Fair soft GDMA & \textbf{0.876} & \textbf{0.822} & \textbf{0.896} & 0.976 & 0.969 & 0.908 \\
\bottomrule
\end{tabular}
\end{table}
